\Cref{tab:phase2} summarises the ablations. The reference-gap and noise conditions use one detector (CNN, seed 0), so they have less statistical power than the main comparison ($n\approx45$ attacks on HAI).

\paragraph{Reference gap $\Delta_{\mathrm{ref}}$} The incident-local reference is taken $\Delta_{\mathrm{ref}}$ seconds before the alarm.
\begin{itemize}
\item \emph{PCL-Sim.} \LX{} is insensitive to the gap. At the alarm it reaches \ptwosimrefALXsingleZerohrOne, \ptwosimrefBLXsingleZerohrOne, \ptwosimrefCLXsingleZerohrOne{} and \ptwosimrefDLXsingleZerohrOne{} for 32, 60, 120 and 300\,s. With \TEA{} at 60\,s it reaches \ptwosimrefALXteaSixZerohrOne--\ptwosimrefDLXteaSixZerohrOne.
\item \emph{HAI.} The default of 60\,s is the best of the four gaps at the alarm (\ptwohairefBLXsingleZerohrOne). A 32\,s gap gives \ptwohairefALXsingleZerohrOne{} and a 300\,s gap \ptwohairefDLXsingleZerohrOne. With \TEA{} at 60\,s the range is \ptwohairefDLXteaSixZerohrOne--\ptwohairefBLXteaSixZerohrOne.
\end{itemize}
A short gap risks a reference that already contains the start of the manipulation. A long gap risks one from a different operating point. On HAI the differences correspond to a few attacks out of about 45, and we did not tune the gap on test data. As expected, RE and \LX-bg do not depend on the reference and are unchanged. IG and KernelSHAP use the same reference as their baseline and vary similarly to \LX.

\paragraph{Sensor noise (RQ4)}
We add i.i.d.\ Gaussian noise with standard deviation $\sigma_n\in\{0.1,0.3\}$ (in standardised units) to the test data.
\begin{itemize}
\item \emph{Clean threshold.} If the residual scale $\sigma$ and the threshold $\theta$ are kept from clean validation data, the detector alarms almost continuously. Precision falls to \ptwohainoiseAcleandetprecision{} on HAI and \ptwosimnoiseAcleandetprecision{} on PCL-Sim, and every attack is ``detected'' by an alarm that is already running at its onset. Localization numbers in this setting measure a saturated detector rather than the explainer, so we do not interpret them. They are included in the released \texttt{phase2.csv}.
\item \emph{Recalibrated threshold.} We therefore also recalibrate $\sigma$ and $\theta$ on validation data with the same noise level, keeping the model weights frozen, as an operator would after commissioning noisier sensors.
\end{itemize}

\noindent With recalibration, the two benchmarks respond very differently.
\begin{itemize}
\item \emph{HAI.} Event recall and delay are nearly unchanged: event recall \ptwohainoiseAdeteventrecall{} and \ptwohainoiseBdeteventrecall, against \ptwohairefBdeteventrecall{} without noise, with median delays of \ptwohainoiseAdetdelaymedians{} and \ptwohainoiseBdetdelaymedians\,s against \ptwohairefBdetdelaymedians\,s. Localization at the alarm degrades for all methods. \LX{} reaches \ptwohainoiseALXsingleZerohrOne{} and \ptwohainoiseBLXsingleZerohrOne{} against \ptwohairefBLXsingleZerohrOne{} without noise; IG reaches \ptwohainoiseAIGsingleZerohrOne{} and \ptwohainoiseBIGsingleZerohrOne. At 60\,s with \TEA, \LX{} reaches \ptwohainoiseALXteaSixZerohrOne{} and \ptwohainoiseBLXteaSixZerohrOne. Against IG and KernelSHAP-1024 with \TEA, \LX{} is numerically lower but not significantly different (Holm $p\ge0.61$, $n=\ptwotesthainoiseAn$ and \ptwotesthainoiseBn).
\item \emph{PCL-Sim.} Noise mainly delays detection. The median delay grows from \ptwosimrefBdetdelaymedians\,s to \ptwosimnoiseAdetdelaymedians\,s and \ptwosimnoiseBdetdelaymedians\,s, presumably because slowly developing manipulations must grow larger before they exceed the wider noise floor. By then they are easy to localize: all methods reach 0.91--0.95 at 60\,s with \TEA, and \LX+\TEA{} does not differ significantly from any of them (Holm $p\ge0.08$, $n=\ptwotestsimnoiseAn$ and \ptwotestsimnoiseBn). The cost appears in the deadline: a localization delivered minutes later is of little use against fast hazards.
\end{itemize}
H5a is therefore \emph{not supported}. \LX+\TEA{} does not degrade less than the baselines under sensor noise. Together with the lower perturbation stability of single \LX{} explanations (\cref{tab:quality}), this suggests that exact loop attributions inherit the noise sensitivity of the detector score. Temporal accumulation still helps under noise (e.g.\ HAI, $\sigma_n=0.1$: \LX{} \ptwohainoiseALXsingleSixZerohrOne{} single-shot vs.\ \ptwohainoiseALXteaSixZerohrOne{} with \TEA{} at 60\,s), but not more for \LX{} than for other methods. Under realistic noise the dominant effect is on detection timeliness, consistent with the deadline decomposition in \cref{sec:results}.


\paragraph{Decoy manipulation (RQ5)} In the decoy suite, every attack is accompanied by a small bias on the PV of an innocent loop. The decoy rarely captures the top rank for any method:
\begin{itemize}
\item at the alarm, \ptwosimdecoyLXsingleZerodecoytopOne{} of attacks for \LX{} and \ptwosimdecoyREsingleZerodecoytopOne{} for RE;
\item at 60\,s with \TEA, \ptwosimdecoyLXteaSixZerodecoytopOne{} for \LX{} and \ptwosimdecoyIGteaSixZerodecoytopOne--\ptwosimdecoyKSOneZeroTwoFourteaSixZerodecoytopOne{} for the other methods.
\end{itemize}
Single-shot explanations 60\,s after the alarm are the most vulnerable. Single-shot \LX{} then ranks the decoy first in \ptwosimdecoyLXsingleSixZerodecoytopOne{} of attacks, because the controllers have partly compensated for the main attack by then. None of the paired differences in decoy capture between \LX+\TEA{} and the other methods with \TEA{} is significant (all Holm $p\ge0.23$, $n=\ptwotestsimdecoyn$ (attack, detector) units). Top-1 accuracy with the decoy is close to that without it (\cref{tab:phase2}). H5b is therefore \emph{not supported}: a non-adaptive decoy of this size is ineffective against all methods, so it cannot discriminate between them. An adaptive adversary remains untested (\cref{sec:limitations}).
